\clearpage
\VerdierPhysicalPageAnchor{0193}%
\begingroup\small\setlength{\parskip}{0pt}\setlength{\emergencystretch}{1.5em}
\setlength{\abovedisplayskip}{4pt plus 2pt minus 2pt}%
\setlength{\belowdisplayskip}{4pt plus 2pt minus 2pt}%

\VerdierSemanticAnchor{ch3d-2-2-6-proof-cont}
(2.1.1)。态射 $p_0$ 是高度为 $1$ 的初等态射。复形 $U_0^\bullet$ 具有形式：

\VerdierSemanticAnchor{ch3d-2-2-6-u0-complex}
\[
\begin{aligned}
U_0^\bullet={}&0\longrightarrow0\longrightarrow\cdots\longrightarrow0
\longrightarrow L_n\longrightarrow L_{n-1}\longrightarrow\cdots\\
&\longrightarrow L_1\longrightarrow Z_0\longrightarrow L
\longrightarrow0\longrightarrow0\longrightarrow\cdots.
\end{aligned}
\]
由性质 $(E_4)$，对象 $Z_0$ 属于 $M'$。因此存在满态射
$u_0\colon L'_0\longrightarrow Z_0$，其中 $L'_0\in M$。如 (2.1.1) 那样，令
$U_1^\bullet=U_0^\bullet(u_0,0)$ 以及
$p_1=p(u_0,0)\colon U_1^\bullet\longrightarrow U_0^\bullet$。

\VerdierSemanticAnchor{ch3d-2-2-6-iterated-elementary-morphisms}
依此继续，归纳地得到一列态射
\[
p_i\colon U_i^\bullet\longrightarrow U_{i-1}^\bullet,
\qquad 1\leq i\leq n-1,
\]
\[
p_0\colon U_0^\bullet\longrightarrow V^\bullet,
\]
其中 $p_i$ 是高度为 $-i+1$ 的初等态射，而 $U_i^\bullet$ 具有形式：

\VerdierSemanticAnchor{ch3d-2-2-6-ui-complex}
\[
\begin{aligned}
U_i^\bullet={}&0\longrightarrow0\longrightarrow\cdots
\longrightarrow L_n\longrightarrow L_{n-1}\longrightarrow\cdots
\longrightarrow L_{i+1}\longrightarrow Z_i\\
&\longrightarrow L'_{i-1}\longrightarrow\cdots\longrightarrow L'_0
\longrightarrow L\longrightarrow0\longrightarrow\cdots ,
\end{aligned}
\]
其中 $L'_j\in M$，$0\leq j\leq i-1$。因此，将这些态射 $p_i$ 复合起来便得到既为满态射又为拟同构的态射

\VerdierSemanticAnchor{ch3d-2-2-6-composite}
\[
p\colon U_{n-1}^\bullet\longrightarrow V^\bullet.
\]
现在注意，在复形

\VerdierSemanticAnchor{ch3d-2-2-6-unminus1-complex}
\[
\begin{aligned}
U_{n-1}^\bullet={}&0\longrightarrow0\longrightarrow\cdots
\longrightarrow L_n\longrightarrow Z_{n-1}\longrightarrow L'_{n-2}
\longrightarrow\cdots\\
&\longrightarrow L'_0\longrightarrow L\longrightarrow0
\longrightarrow\cdots ,
\end{aligned}
\]
中，对象 $Z_{n-1}$ 由性质 $(E_4)$ 表为 $M$ 中两个对象的扩张，故它属于
$M$（由性质 $(E_5)$）。令

\VerdierSemanticAnchor{ch3d-2-2-6-l-complexes}
\[
L^\bullet=0\longrightarrow0\longrightarrow L_n\longrightarrow L_{n-1}
\longrightarrow\cdots\longrightarrow L_0\longrightarrow0
\longrightarrow0\longrightarrow\cdots ,
\]
\[
\begin{aligned}
{L'}^\bullet={}&0\longrightarrow0\longrightarrow L_n
\longrightarrow Z_{n-1}\longrightarrow L'_{n-2}
\longrightarrow\cdots\\
&\longrightarrow L'_0\longrightarrow0\longrightarrow0
\longrightarrow\cdots ,
\end{aligned}
\]
并令 $m\colon {L'}^\bullet\longrightarrow L^\bullet$ 为 $p$ 诱导的态射。
$L^\bullet$（分别地，${L'}^\bullet$）唯一非零的上同调对象是 $Y$，位于次数
$0$（分别地，是 $L$，位于次数 $0$），而 $m$ 在这些对象上诱导的态射是 $u$。
考虑态射 $m$ 的映射锥 $c(m)$：

\VerdierSemanticAnchor{ch3d-2-2-6-cone}
\[
\begin{aligned}
c(m)={}&\cdots\longrightarrow0\longrightarrow L_n
\longrightarrow L_n\oplus Z_{n-1}
\longrightarrow L_{n-1}\oplus L'_{n-2}\longrightarrow\cdots\\
&\longrightarrow L_1\oplus L'_0\longrightarrow L_0
\longrightarrow0\longrightarrow\cdots.
\end{aligned}
\]
\endgroup
